\uuid{Arnn}
\titre{Changements de variables/primitives usuelles}
\auteur{hanani ; mzari}
\organisation{exo7}
\licence{CC BY-NC-SA 4.0}
\datecreate{2020-01-20}
\niveau{L1}
\module{Analyse}
\chapitre{Calcul d'intégrales}
\sousChapitre{Primitives diverses}
\difficulte{2}
\theme{Primitives, changements de variables, primitives usuelles}
\exo7qcmid{lille-2.primitives.014}

\contenu{
\begin{question}
\qclassification{127.11}
\qtags{motcle=changements de variables/primitives usuelles}
Parmi les assertions suivantes, cocher celles qui sont vraies :
\begin{answers}
\bad{Les primitives de $16x^3(x^4+1)^3$ sur $\Rr$ sont données par $\displaystyle F(x)=x^4(x^4+1)^4+k$, $k\in \Rr$.}
\good{Les primitives de $16x^3(x^4+1)^3$ sur $\Rr$ sont données par $\displaystyle F(x)=(x^4+1)^4+k$, $k\in \Rr$.}
\good{Les primitives de $\displaystyle \frac{4x^3}{\sqrt{x^4+1}}$ sur $\Rr$ sont données par $\displaystyle F(x)=2\sqrt{x^4+1}+k$, $k\in \Rr$.}
\bad{Les primitives de $\displaystyle \frac{4x^3}{\sqrt{x^4+1}}$ sur $\Rr$ sont données par $\displaystyle F(x)=\frac{x^4}{\sqrt{x^4+1}}+k$, $k\in \Rr$.}
\end{answers}
\begin{explanations}
Avec $u=x^4+1$, on a : $\mathrm{d}u=4x^3\, \mathrm{d}x$. D'où
$$\int 16x^3(x^4+1)^3\mathrm{d}x=\int 4u^3\, \mathrm{d}u=u^4+k=(x^4+1)^4+k,\; k\in \Rr,$$
et
$$\int \frac{4x^3\, \mathrm{d}x}{\sqrt{x^4+1}}=\int \frac{\mathrm{d}u}{\sqrt{u}}=2\sqrt{u}+k=2\sqrt{x^4+1}+k,\; k\in \Rr.$$
\end{explanations}
\end{question}
}
